\chapter{Teoria cinetica dei gas}

L'equazione di stato dei gas perfetti $PV = nRT$ descrive il gas dal punto di vista macroscopico, cioè come un oggetto di dimensioni paragonabili alle nostre, fatto di litri e di atmosfere, senza dire nulla su che cosa il gas sia.
La teoria cinetica\mkeyword{teoria cinetica} colma questa lacuna: assume che il gas sia costituito da molecole in moto e ricava l'equazione di stato dalla meccanica di quelle molecole.
Il risultato è notevole non tanto perché riproduce una legge già nota, quanto perché nel farlo attribuisce un significato meccanico a due grandezze, la pressione e la temperatura, che nella termodinamica macroscopica sono primitive.

\section{Il gas dal punto di vista microscopico}

Una mole contiene $N_A = 6,022\cdot 10^{23}$ molecole, e la prima osservazione da fare è che in un gas queste molecole occupano una frazione trascurabile dello spazio disponibile.
Una mole di acqua gassosa a $1 \ atm$ e $300 \ K$ occupa
\begin{equation}
    V_m = \frac{RT}{P} = \frac{0,0821 \cdot 300}{1} \simeq 24,6 \ L,
    \label{eq:volumemolare}
\end{equation}
mentre la stessa mole allo stato liquido pesa $18 \ g$ e, avendo l'acqua densità unitaria, occupa $18 \ mL$.
Il rapporto fra i due volumi vale
\begin{equation}
    \frac{18 \ mL}{24,6 \ L} = 7,3\cdot 10^{-4},
    \label{eq:rapportovolumi}
\end{equation}
e poiché nel liquido le molecole sono già a contatto, la \eqref{eq:rapportovolumi} è una sovrastima del volume proprio delle molecole: almeno il $99{,}93\%$ del volume di un gas è spazio vuoto.
Se le molecole occupano così poco spazio e il gas riempie ugualmente tutto il recipiente, l'unica spiegazione possibile è che esse si muovano.

\section{Interpretazione microscopica della pressione}

La pressione è la conseguenza degli urti delle molecole contro le pareti del recipiente.
Consideriamo una molecola di massa $m$ che urta elasticamente una parete perpendicolare all'asse $x$: prima dell'urto la sua velocità è $\mathbf{v} = (v_x, v_y, v_z)$, dopo l'urto è $(-v_x, v_y, v_z)$, e la quantità di moto trasferita alla parete vale

\begin{equation}
    \Delta p = 2mv_x .
    \label{eq:qdmurto}
\end{equation}

Se il recipiente è un cubo di lato $L$, fra due urti successivi sulla stessa parete la molecola percorre un tratto $2L$ lungo $x$, cosicché l'intervallo fra due urti è $\Delta t = 2L/v_x$.
La forza media esercitata da quella singola molecola su quella singola parete è dunque

\begin{equation}
    F = \frac{\Delta p}{\Delta t} = \frac{2mv_x}{2L/v_x} = \frac{mv_x^2}{L}.
    \label{eq:forzasingola}
\end{equation}
Sommando su tutte le $N$ molecole e introducendo la media $\langle v_x^2\rangle$, poiché le tre direzioni dello spazio sono equivalenti e vale quindi $\langle v_x^2\rangle = \langle v^2\rangle/3$, si ottiene per la forza totale $Nm\langle v^2\rangle/(3L)$.
Dividendo per la superficie $L^2$ della parete e riconoscendo $V = L^3$ si arriva al risultato centrale della teoria cinetica,

\begin{equation}
    PV = \frac{1}{3}Nm\langle v^2\rangle .
    \label{eq:cineticabase}
\end{equation}

\subsection{La distribuzione di Maxwell-Boltzmann}

Non tutte le molecole si muovono alla stessa velocità, e la distribuzione dei moduli è quella di Maxwell-Boltzmann\mkeyword{distribuzione di Maxwell-Boltzmann}

\begin{equation}
    f(v) = 4\pi\left(\frac{m}{2\pi k_BT}\right)^{3/2} v^2 \, e^{-mv^2/2k_BT},
    \label{eq:maxwell}
\end{equation}

dove $f(v)\,dv$ è la frazione di molecole con modulo della velocità compreso fra $v$ e $v + dv$.
La struttura della \eqref{eq:maxwell} spiega da sola la forma del grafico: il fattore gaussiano $e^{-mv^2/2k_BT}$ penalizza le velocità alte, mentre il fattore $v^2$, che proviene dall'elemento di volume in coordinate sferiche nello spazio delle velocità, penalizza quelle basse.

\begin{figure}[htbp]
    \centering
    \begin{tikzpicture}[x=3.6cm, y=2.9cm, >={Latex[length=2mm]}]
        % --- macro riutilizzabili -----------------------------------------
        \def\mb{2.71828*pow(\x,2)*exp(-pow(\x,2))}
        \def\vp{1}
        \def\vmed{1.1284}
        \def\vrms{1.2247}
        \def\marca#1#2#3{\draw[dashed] (#1,#3) -- (#1,#2);}
        % --- assi ---------------------------------------------------------
        \draw[->] (-0.05,0) -- (2.55,0) node[below left] {velocità $v$};
        \draw[->] (0,-0.05) -- (0,1.20) node[right] {$f(v)$};
        % --- curva --------------------------------------------------------
        \draw[domain=0:2.5, samples=200, smooth, thick] plot (\x, {\mb});
        % --- velocita' caratteristiche ------------------------------------
        \marca{\vp}{1.000}{-0.06}
        \marca{\vmed}{0.969}{-0.20}
        \marca{\vrms}{0.910}{-0.34}
        \node[font=\footnotesize, anchor=north] at (\vp,-0.06) {$v_p$};
        \node[font=\footnotesize, anchor=north] at (\vmed,-0.20) {$\langle v\rangle$};
        \node[font=\footnotesize, anchor=north] at (\vrms,-0.34) {$v_{rms}$};
    \end{tikzpicture}
    \caption{Distribuzione di Maxwell-Boltzmann dei moduli della velocità.}
    \label{fig:maxwell}
\end{figure}

Il prodotto di un fattore crescente e di uno decrescente produce una curva asimmetrica, ripida a sinistra e con una coda lunga a destra, la cui area resta unitaria a ogni temperatura.
Al crescere di $T$, o al diminuire della massa, la curva si abbassa e si allarga verso destra, ma il suo integrale non cambia.

L'asimmetria ha una conseguenza che conviene rendere quantitativa.
Dalla \eqref{eq:maxwell} si ricavano tre velocità caratteristiche distinte,

\begin{equation}
    v_p = \sqrt{\frac{2k_BT}{m}},
    \qquad
    \langle v\rangle = \sqrt{\frac{8k_BT}{\pi m}},
    \qquad
    v_{rms} = \sqrt{\langle v^2\rangle} = \sqrt{\frac{3k_BT}{m}},
    \label{eq:velocitacaratteristiche}
\end{equation}

cioè la velocità più probabile, che corrisponde al massimo della curva, la velocità media e la velocità quadratica media.

\section{Temperatura ed energia cinetica}

Riscrivendo la \eqref{eq:cineticabase} in termini dell'energia cinetica media di traslazione $\varepsilon_{tr} = \tfrac{1}{2}m\langle v^2\rangle$ si ottiene

\begin{equation}
    PV = \frac{2}{3}N\varepsilon_{tr},
    \label{eq:pvenergia}
\end{equation}

e il confronto con l'equazione di stato scritta nella forma $PV = Nk_BT$ fornisce immediatamente

\begin{equation}
    \varepsilon_{tr} = \frac{3}{2}k_BT,
    \qquad
    E_{tr,m} = \frac{3}{2}RT .
    \label{eq:equipartizione}
\end{equation}

La temperatura assoluta è dunque, a meno di un fattore, l'energia cinetica media di traslazione di una molecola, e questo è il contenuto fisico del principio di equipartizione\mkeyword{equipartizione}: ciascuno dei tre gradi di libertà traslazionali contribuisce $\tfrac{1}{2}k_BT$.\\

Va aggiunto che la \eqref{eq:equipartizione} riguarda la sola energia di traslazione.
Per un gas monoatomico come l'elio essa esaurisce l'energia interna, e da qui segue il calore specifico $C_V = \tfrac{3}{2}R$; per una molecola biatomica si aggiungono due gradi di libertà rotazionali, e per una poliatomica non lineare tre, con i corrispondenti contributi al calore specifico.
L'equazione di stato non ne risente, perché la pressione dipende solo dal moto del centro di massa, ed è questa la ragione per cui $PV = nRT$ vale identica per gas mono e poliatomici mentre i loro calori specifici differiscono.

\subsection{Effusione e legge di Graham}

Poiché a parità di temperatura l'energia cinetica media è la stessa per tutti i gas, dalla condizione
\begin{equation}
    \frac{1}{2}m_1\langle v_1^2\rangle = \frac{1}{2}m_2\langle v_2^2\rangle
    \qquad\Longrightarrow\qquad
    \frac{\langle v_1^2\rangle}{\langle v_2^2\rangle} = \frac{m_2}{m_1}
    \label{eq:graham}
\end{equation}
segue che le velocità stanno fra loro come l'inverso della radice delle masse.
Il flusso di molecole che attraversa un piccolo foro è proporzionale a $\langle v\rangle$, e dunque la velocità con cui due gas fuoriescono da uno stesso foro sta nel rapporto $\sqrt{M_2/M_1}$: è la legge di Graham, che consente di separare gas diversi in base alla massa molecolare e che ha trovato la sua applicazione storicamente più rilevante nell'arricchimento isotopico dell'uranio, dove la separazione fra $^{235}\mathrm{UF_6}$ e $^{238}\mathrm{UF_6}$ sfrutta un rapporto di masse di appena $1,0043$, con un fattore di separazione per stadio di $1,0043^{1/2} \simeq 1,0021$ che rende necessarie migliaia di stadi in cascata.

\newpage

\section{Equazione di Van der Waals}

Il modello di gas perfetto poggia su due ipotesi: che il volume proprio delle molecole sia trascurabile rispetto a quello del recipiente, cioè che le molecole siano puntiformi, e che fra di esse non agiscano forze a distanza.
Entrambe valgono rigorosamente solo nel limite di basse pressioni e grandi volumi molari, dove le molecole sono in media lontane fra loro.

Il fallimento più vistoso è però qualitativo prima ancora che quantitativo: l'equazione $PV = nRT$ non ammette alcuna transizione di fase, e prevede che un gas resti gas a qualunque pressione.
Occorre dunque un'equazione di stato che contenga, sia pure in forma rudimentale, le due proprietà che il gas perfetto ignora, cioè il volume delle molecole e le forze intermolecolari studiate nel capitolo precedente.

Conviene passare a grandezze intensive, introducendo il volume molare $V_m = V/n$ e la densità numerale $\rho_n = n/V = 1/V_m$, e scrivere l'equazione di stato nella forma $PV_m = RT$.

\subsection{Il covolume}

La prima correzione riguarda il volume.
Le molecole hanno dimensioni finite, quindi il volume effettivamente disponibile per il loro moto non è il volume del recipiente ma quello diminuito di un termine $b$, detto covolume\mkeyword{covolume}, con le stesse dimensioni di un volume molare:

\begin{equation}
    V_m^{gp} = V_m - b .
    \label{eq:covolume}
\end{equation}

A differenza di $R$, il covolume non è una costante universale, ma dipende dal gas, il che è del resto ovvio se ha il significato che gli si attribuisce.

\subsection{Il termine attrattivo}

La seconda correzione riguarda la pressione.
Una molecola nel corpo del gas è circondata da vicini in tutte le direzioni e la risultante delle forze attrattive che subisce è nulla in media; una molecola che sta per urtare la parete, invece, ha vicini soltanto dalla parte del gas, e viene quindi frenata verso l'interno.
Arriva alla parete con quantità di moto minore e la pressione effettivamente misurata è minore di quella che si osserverebbe in assenza di interazioni:

\begin{equation}
    P^{gp} = P + \frac{a}{V_m^2}.
    \label{eq:correzionepressione}
\end{equation}

La dipendenza da $V_m^{-2}$ segue da un argomento di densità di coppie.
L'entità della frenata è proporzionale al numero di molecole che tirano indietro quella incidente, cioè a $\rho_n = 1/V_m$, e il numero di molecole che arrivano alla parete nell'unità di tempo è a sua volta proporzionale a $\rho_n$: il prodotto delle due densità dà $\rho_n^2 = 1/V_m^2$, e $a$ è la costante di proporzionalità.\\

Mettendo insieme le due correzioni si ottiene l'equazione di van der Waals\mkeyword{equazione di van der Waals}

\begin{equation}
    \left(P + \frac{a}{V_m^2}\right)(V_m - b) = RT .
    \label{eq:vanderwaals}
\end{equation}

\section{Isoterme e punto critico}

Esplicitando la \eqref{eq:vanderwaals} rispetto alla pressione si ottiene

\begin{equation}
    P = \frac{RT}{V_m - b} - \frac{a}{V_m^2},
    \label{eq:vdwesplicita}
\end{equation}

espressione nella quale il primo termine, repulsivo, domina a grandi volumi molari, cioè a basse pressioni e alte temperature, mentre il secondo, attrattivo, prende il sopravvento a volumi molari piccoli.\\

Studiando la \eqref{eq:vdwesplicita} a $T$ fissata si trovano tre famiglie di isoterme.

Al di sopra di una temperatura caratteristica $T_c$ le isoterme sono monotone decrescenti e il gas si comporta qualitativamente come un gas perfetto, tanto più fedelmente quanto più ci si allontana da $T_c$. L'isoterma prende il nome di isoterma supercritica\mkeyword{isoterma supercritica}.

Al di sotto di $T_c$ le isoterme presentano un massimo e un minimo relativi. L'isoterma prende il nome di isoterma subcritica\mkeyword{isoterma subcritica}.

Esattamente a $T = T_c$ massimo e minimo si fondono in un flesso a tangente orizzontale, che definisce il punto critico\mkeyword{punto critico}. L'isoterma prende il nome di isoterma critica\mkeyword{isoterma critica}.

\begin{figure}[htbp]
    \centering
    \begin{tikzpicture}[x=1.6cm, y=3.0cm, >={Latex[length=2mm]}]
        % --- macro riutilizzabili -----------------------------------------
        \def\iso#1{8*#1/(3*\x-1)-3/(\x*\x)}
        \def\isoterma#1#2{\draw[domain=#2:3.60, samples=220, smooth, thick] plot (\x, {\iso{#1}});}
        % --- assi ---------------------------------------------------------
        \draw[->] (0,0) -- (4.35,0) node[below left] {$V_m$};
        \draw[->] (0,0) -- (0,2.15) node[below right] {$P$};
        % --- isoterme -----------------------------------------------------
        \isoterma{1.15}{0.700}
        \isoterma{1.00}{0.573}
        \isoterma{0.85}{0.500}
        % --- etichette delle isoterme -------------------------------------
        \node[font=\scriptsize, anchor=west] at (3.66,0.746) {$T > T_c$};
        \node[font=\scriptsize, anchor=west] at (3.66,0.624) {$T = T_c$};
        \node[font=\scriptsize, anchor=west] at (3.66,0.502) {$T < T_c$};
        % --- punto critico ------------------------------------------------
        \fill (1,1) circle (1.6pt);
        \node[font=\scriptsize, anchor=south east] at (1.90,1.00) {$(V_{mc},P_c)$};
        % --- segmento di Maxwell ------------------------------------------
        \draw[dashed] (0.553,0.5045) -- (3.128,0.5045);
    \end{tikzpicture}
    \caption{Isoterme di van der Waals.}
    \label{fig:isotermevdw}
\end{figure}

L'isoterma critica separa due regioni del piano $P$--$V_m$ con comportamenti qualitativamente diversi.
Al di sopra di $T_c$ il gas non può essere liquefatto per sola compressione, per quanto grande sia la pressione applicata; al di sotto la compressione produce liquefazione, e in questo regime è appropriato chiamare la sostanza vapore anziché gas.

\subsection{Isoterma supercritica}

Per $T > T_c$ la pressione decresce monotonamente al crescere del volume molare e l'isoterma non presenta punti stazionari.
Comprimendo il fluido lungo una di queste curve la pressione sale indefinitamente senza che compaia alcuna discontinuità, e la sostanza non può essere liquefatta per sola compressione, per quanto grande sia la pressione applicata.
Quanto più ci si allontana da $T_c$ tanto più il termine attrattivo diventa irrilevante rispetto a $RT/(V_m-b)$, e la curva tende all'iperbole del gas perfetto.

\subsection{Isoterma critica}

A $T = T_c$ il massimo e il minimo relativi si fondono in un unico punto di flesso a tangente orizzontale, definito dall'annullamento simultaneo delle prime due derivate,
\begin{equation}
    \left(\frac{\partial P}{\partial V_m}\right)_{T_c} = 0,
    \qquad
    \left(\frac{\partial^2 P}{\partial V_m^2}\right)_{T_c} = 0 .
    \label{eq:flesso}
\end{equation}
Imponendo le \eqref{eq:flesso} alla \eqref{eq:vdwesplicita} ed eliminando la temperatura fra le due condizioni si ricava $V_m - b = \tfrac{2}{3}V_m$, da cui seguono le coordinate del punto critico in funzione delle sole costanti del gas:
\begin{equation}
    V_{mc} = 3b,
    \qquad
    T_c = \frac{8a}{27Rb},
    \qquad
    P_c = \frac{a}{27b^2}.
    \label{eq:costanticritiche}
\end{equation}
Le \eqref{eq:costanticritiche} rendono conto di due fatti familiari: l'elio, con $T_c = 5{,}2 \ K$, va raffreddato a temperature estreme prima di poter essere liquefatto, mentre l'anidride carbonica, con $T_c = 304 \ K$, liquefa già a temperatura ambiente sotto circa $70 \ atm$, ed è per questo che le bombole di $\mathrm{CO_2}$ contengono liquido e quelle di azoto gas compresso.

\subsection{Isoterma subcritica}

Per $T < T_c$ l'isoterma presenta un minimo e un massimo relativi, e nel tratto che li congiunge si avrebbe $(\partial P/\partial V_m)_T > 0$, cioè una pressione che cresce all'aumentare del volume.
Uno stato simile è meccanicamente instabile e non è osservabile: una fluttuazione che comprimesse localmente il fluido ne abbasserebbe la pressione, favorendo un'ulteriore compressione in un processo che si autoamplifica.
L'ansa va quindi sostituita con un segmento orizzontale, posto in modo che le due aree che esso racchiude con la curva siano uguali, condizione nota come costruzione di Maxwell\mkeyword{costruzione di Maxwell} ed equivalente all'uguaglianza dei potenziali chimici delle due fasi agli estremi del segmento.
Percorrendo l'isoterma in compressione si distinguono allora tre tratti: a destra del segmento si comprime il vapore, con andamento simile a quello supercritico; lungo il segmento il vapore liquefa progressivamente a pressione costante e le due fasi coesistono; a sinistra la curva diventa ripidissima, il che è la traduzione grafica del fatto che un liquido è quasi incomprimibile.